\subsection{Types of morphisms}

\subsubsection{Open embeddings, finite morphisms}

Recall that a morphism of locally ringed spaces $(X,\O_X)\to(Y,\O_Y)$ is a topological map $\phi\colon X\to Y$
equipped with a morphism of sheaves $\phi^\sharp\colon\O_X\to\phi_*\O_Y$ such that the associated map
of stalks $\phi_p\colon\O_{X,\phi p}\to\O_{Y,p}$ maps the maximal ideal $\m_{X,\phi p}$ into $\m_{Y,p}$
for all $p\in Y$.
A morphism of schemes is then a morphism of locally ringed spaces.
It is clear that an isomorphism of ringed spaces is an isomorphism of locally ringed spaces.

Furthermore recall the mutual right adjunction
$$ \hom_{\LRS}(X,\spec A) \cong \hom_\Ring(A,\Gamma(X,\O_X)) $$
which was proven in \refmath[theorem]{specglobsectionadj}.

The affine scheme $\spec\bZ$ is terminal in the category of affine schemes (because $\bZ$ is the initial ring),
and it is thus terminal in the category of schemes by taking an affine cover.

\bdefn

    Let $S$ be a scheme.
    The category of {\emphcolor $S$-schemes} is the over category $\Sch_S$.
    That is, it is the category whose objects are schemes $X$ with a morphism $X\to S$ (called the
    {\emphcolor structure morphism}), and morphisms
    are commutative triangles

    \diagram{
        X&&Y\cr
        &S
    }{
        \diagarrow{from={1,1}, to={1,3}}
        \diagarrow{from={1,1}, to={2,2}}
        \diagarrow{from={1,3}, to={2,2}}
    }

\edefn

For a ring $A$, we may define an $A$-scheme to be a scheme whose structure sheaf is valued in the category
of $A$-algebras.
A $\spec A$-scheme determines an $A$-algebra, as the structure morphism gives a global section
$A\to\Gamma(X,\O_X)$ and all sections are $\Gamma(X,\O_X)$-algebras.
The converse follows by \refmath[theorem]{specglobsectionadj}.
So an $A$-scheme is the same as a $\spec A$-scheme.

\bdefn

    A class of morphisms of schemes is called {\emphcolor reasonable} if it contains all isomorphisms
    and satisfies the following properties:
    \benum
        \item The class is closed under composition,
        \item The class is preserved by base change (pullback): if $X\to Y$ is in the class, then for any
        $Y'\to Y$, the induced map $X\times_YY'\to Y'$ is also in the class.
        (This assumes the pullback exists, and we will see that pullbacks of schemes always exist.)
        \item The class is {\emphcolor local on target}.
        That is to say, (a) if $\pi\colon X\to Y$ is in the class, then for any open $V\subseteq Y$, the
        restriction $\pi^{-1}V\to V$ is in the class.
        And (b) if there is an open cover $Y=\bigcup_iV_i$ such that $\pi^{-1}V_i\to V_i$ is in the class for
        all $i$ then $\pi$ is in the class.
    \eenum

\edefn

Note that if a property has (1) and (2) then it is closed under products.
That is to say if $X\to Y$ and $X'\to Y'$ are morphisms of $S$-schemes with the property then
$X\times_SX'\to Y\times_SY'$ has the property.
Indeed we note that $X\times_SX'\to X'\times_SY$ has the property as it is the morphism
$$ X\times_SX' \cong X\times_Y(Y\times_SX') \to Y\times_SX' \cong X'\times_SY . $$
Similarly $X'\times_SY\to X'\times_SY'$ has the property, and thus $X\times_SX'\to Y\times_SY'$ does.

\bexam

    Clearly the class of isomorphisms is reasonable.

\eexam

\bdefn

    A morphism $X\to Y$ is an {\emphcolor open embedding} (or {\emphcolor open immersion}) if it factors
    as
    $$ (X,\O_X) \lisoto (U,\O_Y|_U) \inj (Y,\O_Y) $$
    for $U\subseteq Y$ open.
    That is $X$ is isomorphic to an open subscheme of $Y$.

\edefn

Open embeddings clearly satisfy (1) and (3) of being reasonable.
Now if $\iota\colon U\to Z$ is an open embedding and $\rho\colon Y\to Z$ is an arbitrary morphism, then
$U\times_ZY=(\rho^{-1}U,\O_Y|_{\rho^{-1}U})$.
This is easy enough to check, and this is clearly an open subscheme of $Y$.
Thus the class of open embeddings is reasonable.

If $X\subseteq Y$ is an open subscheme and $Y$ is locally Noetherian, then so too is $X$ (clearly, as an
affine open subscheme of $X$ is an affine open subscheme of $Y$).
And if $Y$ is Noetherian so too is $X$ as $Y$ is a Noetherian space and thus $X$ is quasicompact.

\bdefn

    A property of morphisms is {\emphcolor local on the source} when (a) if $\pi\colon X\to Y$
    has the property and $U\subseteq X$ is an open subscheme, then $\pi|_U\colon U\to Y$ has the property.
    And (b) if $X=\bigcup_iU_i$ is an open cover where $\pi|_{U_i}\colon U_i\to Y$ has the property then
    so does $\pi$.

\edefn

\bdefn

    Let $\pi\colon X\to Y$ be a morphism of topological spaces.
    It is {\emphcolor quasicompact} if for every quasicompact open $U\subseteq Y$, $\pi^{-1}U\subseteq X$
    is quasicompact.
    And it is {\emphcolor quasiseparated} if for every quasiseparated open $U\subseteq Y$,
    $\pi^{-1}U\subseteq X$ is quasiseparated.

    A morphism of schemes has one of these properties if the underlying topological map has it.

\edefn

For a morphism of schemes it is sufficient to check quasicompactness/quasiseparatedness on affine open subsets.
This is because any open subset of a scheme can be covered with an affine open cover, and if it is
quasicompact then this cover can be taken to be finite, and its preimage is a finite union of quasicompact
spaces and thus quasicompact.
And similarly for quasiseparatedness: a space with an open cover of quasiseparated subspaces is quasiseparated.

The composition of quasicompact or quasiseparated morphisms is quasicompact or quasiseparated respectively.

\bexam

    Clearly if $X$ is a Noetherian space then any map $X\to Y$ is quasicompact.
    In particular if $X$ is a Noetherian scheme (this is not true for $X$ quasicompact!).
    And if $X$ is a quasiseparated space then any map $X\to Y$ is quasiseparated (as all open subsets of
    $X$ are quasiseparated).

\eexam

\bprop

    Let $\pi\colon X\to Y$ be a morphism of schemes.
    \benum
        \item $\pi$ is quasicompact iff there is an open affine cover $Y=\bigcup_iU_i$ such that $\pi^{-1}U_i$ is quasicompact.
        \item $\pi$ is quasiseparated iff there is an open affine cover $Y=\bigcup_iU_i$ such that $\pi^{-1}U_i$ is quasiseparated.
    \eenum

\eprop

\bproof

    Both points are proven with affine communication.
    \qed

\eproof

\bdefn

    A morphism of schemes $\pi\colon X\to Y$ is {\emphcolor affine} if $\pi^{-1}U$ is affine for all $U\subseteq Y$ affine.

\edefn

Because quasicompactness and quasiseparatedness can be checked on affine open subsets of $Y$, affine morphisms are quasicompact and
quasiseparated.

\bprop

    A morphism $\pi\colon X\to Y$ is affine iff there is an open affine cover $Y=\bigcup_iU_i$ such that $\pi^{-1}U_i$ are all affine.

\eprop

\bproof

    We check that the property of an affine $U\subseteq Y$ having $\pi^{-1}U$ being affine satisfies the conditions of the
    \refmath{affcomm}.
    Clearly for $\spec A\subseteq Y$ if $\pi^{-1}\spec A=\spec B$ then $\pi^{-1}\spec A_f=\spec B_{\pi^\sharp f}$.
    And if $(f_1,\dots,f_n)=A$ and $\pi^{-1}\spec A_{f_i}=\spec B_i$.
    Now we may suppose that our map is $\pi\colon X\to\spec A$, i.e.\ the codomain is $\spec A$, and we want to show $X$ is affine.
    So define $B=\Gamma(X,\O_X)$ so that $X\to\spec A$ factors through $X\to\spec B$ (arising from the identity $B\to\Gamma(X,\O_X)$).

    So we have the commutative triangle
    \diagram{
        X&&\spec B\cr
        &\spec A
    }{
        \diagarrow{from={1,1}, to={1,2}, text=\alpha, y distance=.15cm}
        \diagarrow{from={1,1}, to={2,2}, text=\pi, x distance=.15cm}
        \diagarrow{from={1,2}, to={2,2}, text=\beta, x distance=.15cm}
    }
    And we want to show that $\alpha$ is an isomorphism.
    Now $\spec A=\bigcup_i\dof(f_i)$ and $\beta^{-1}\dof(f_i)=\dof(\beta^\sharp f_i)=\spec B_{\beta^\sharp f_i}$.
    And $\pi^{-1}\dof(f_i)=\spec B_i$.
    Furthermore we have that $X=\pi^{-1}\spec B$ is quasicompact and quasiseparated by the above proposition.

    By \refmath{qcqslemm} there is a natural isomorphism $B_{\beta^\sharp f_i}=\Gamma(X,\O_X)_{\beta^\sharp f_i}\cong\Gamma(X_{\beta^\sharp f_i},\O_X)=A_i$.
    And so $\alpha$ induces an isomorphism $\spec B_i\to\spec B_{\beta^\sharp f_i}$.
    Because the $B_i$ cover $\spec B$ we have that $\alpha$ is an isomorphism $X\to\spec B$ as desired.
    \qed

\eproof

\bcoro

    Let $Z\subseteq\spec A$ be closed such that locally $Z$ is cut out by a single equation.
    That is, there is an open affine cover $\spec A=\bigcup_i\spec A_i$ such that $Z\cap\vof_{A_i}(f_i)$ for $f_i\in A_i$.
    Then $Y=\spec A-Z$ is affine.

\ecoro

\bproof

    Under the inclusion $\iota\colon Y\to\spec A$, the preimage of each $\spec A_i\subseteq\spec A$ is $\dof_{A_i}(f_i)=\spec A_{i,f_i}$ which is affine.
    By the above proposition $\iota$ is an affine morphism, and so $Y=\iota^{-1}\spec A$ is affine.
    \qed

\eproof

Recall that a morphism of rings $A\to B$ is {\emphcolor finite} if it endows $B$ with an $A$-algebra structure such that $B$ is finitely generated as an $A$-{\it module}.
And it is of {\emphcolor finite type} if $B$ is a finitely generated $A$-{\it algebra}.

\bdefn

    A morphism of schemes $\pi\colon X\to Y$ is {\emphcolor finite} if it is affine and for $\spec B\subseteq Y$, $\pi^{-1}\spec B=\spec A$
    and $\pi\colon\spec A\to\spec B$ makes $A$ into a finite $B$-algebra.

    It is {\emphcolor locally of finite type} if for every affine open $\spec A\subseteq X$ and $\spec B\subseteq Y$ such that $\pi\spec A\subseteq\spec B$,
    $A$ is a finite type $B$-algebra.
    It is {\emphcolor of finite type} if it is locally of finite type and quasicompact (so the preimage of affine open in $Y$ can be covered by finitely many
    affine opens in $X$).

\edefn

\bprop

    $\pi\colon X\to Y$ is locally of finite type iff $Y=\bigcup_i\spec B_i$ and $\pi^{-1}\spec B_i=\bigcup_j\spec A_{ij}$ where $A_{ij}$ is a finite type $B_i$-algebra.
    And it is of finite type if there are finitely many $A_{ij}$.

\eprop

\bproof

    See homework 2.
    \qed

\eproof

\bprop

    $\pi\colon X\to Y$ is finite iff $Y=\bigcup_i\spec B_i$ and $\pi^{-1}\spec B_i=\spec A_i$ making $A_i$ into a finite $B_i$-algebra.

\eprop

\bproof

    By affine communication, if $\pi^{-1}\spec B=\spec A$ then $\pi^{-1}\spec B_f=\spec A_{\pi^\sharp f}$ and $A_{\pi^\sharp f}$ is a finite $B_f$-algebra.
    Conversely if $B=(f_1,\dots,f_n)$ and $A_{\pi^\sharp f_i}$ are finite $B_{f_i}$-algebras, then $A$ is a finite $B$-algebra.
    \qed

\eproof

\subsubsection{Chevalley's theorem}

\TODO{This}

\subsubsection{Closed embeddings}

\bdefn

    A morphism of schemes $\pi\colon X\to Y$ is a {\emphcolor closed embedding} if it is affine, and for $\spec B\subseteq Y$ and
    $\pi^{-1}\spec B=\spec A$, the induced map $B\to A$ is surjective.
    A {\emphcolor closed subscheme} of a scheme $Y$ is a subset $X\subseteq Y$ which is a scheme and for which the inclusion is a closed
    embedding.

\edefn

Clearly closed embeddings are finite and closed embeddings are closed under composition.
Furthermore by affine locality we may show easily that the property of being a closed embedding is affine-local on the target.

\bremark

    If $\pi\colon X\to Y$ is a closed embedding, then $\pi$ is a topological closed embedding $\pi^\sharp\colon\O_Y\to\pi_*\O_X$ is surjective.
    Indeed, if $Y=\bigcup_i\spec B_i$ is an affine open cover and $\pi^{-1}\spec B_i=\spec A_i$ then the restriction
    $\pi\colon\spec A_i\to\spec B_i$ is a closed embedding (corresponding to $\spec B/I\cong\vof(I)\inj\spec B$).
    Because being a topological closed embedding is local on target, $\pi$ is.

    Thus the stalks of $\pi_*\O_X$ are $(\pi_*\O_X)_{\pi p}\cong\O_{X,p}$ for $p\in X$ and otherwise zero when the point is outside of
    $\pi$'s image.
    Then on $\spec B\subseteq Y$ we have that $\O_Y|_{\spec B}\to(\pi_*\O_X)|_{\spec B}$ is $\spec B\to\spec A$, whose stalks are surjective.
    Thus $\pi^\sharp$'s stalks are surjective, and thus $\pi^\sharp$ is.

    The converse is also true: if $\pi$ is a topological closed embedding and $\pi^\sharp$ is an epimorphism of sheaves, then $\pi$ is a closed
    embedding.

\eremark

\bdefn

    Let $Y$ be a scheme (or a ringed space), then a sub-$\O_Y$-module of $\O_Y$ is called an {\emphcolor ideal sheaf}.
    This is equivalent to its sections being ideals of $\O_Y$'s sections.

\edefn

\blemm

    If $\pi\colon X\to Y$ is affine then $\pi_*\O_X$ is a quasicoherent sheaf on $Y$.

\elemm

\bproof

    By \refmath[lemma]{distinguishedquasicoherent}, it is sufficient to show that for $\spec A\subseteq Y$ and $f\in A$,
    we have $\Gamma(\spec A,\pi_*\O_X)_f\cong\Gamma(\spec A_f,\pi_*\O_X)$.
    By assumption $\pi^{-1}\spec A=\spec B$ for some $B$ and so $\pi^{-1}\spec A_f=\spec B_{\pi^\sharp f}$.
    Thus we have the desired isomorphism.
    \qed

\eproof

In particular if $\pi$ is a closed embedding then $\pi_*\O_X$ is a quasicoherent sheaf on $Y$.
Because quasicoherent sheaves form an Abelian category, the kernel of $\O_Y\to\pi_*\O_X$ is quasicoherent.
Call it $\I_{X/Y}$, and it fits into an exact sequence of quasicoherent sheaves on $Y$
$$ 0 \to \I_{X/Y} \to \O_Y \to \pi_*\O_X \to 0 , $$
called the {\emphcolor closed subscheme exact sequence} for $\pi\colon X\inj Y$.
From such an exact sequence we recover $X$ and $\pi$.
Taking local sections on an affine open $\spec B\subseteq Y$ is exact on quasicoherent sheaves, because quasicoherent sheaves on $\spec B$
are in equivalence with $B$-modules (via global sections and $M\mapsto\tilde M$).
So if $I=\Gamma(\spec B,\I_{X/Y})$ then we have an exact sequence
$$ 0 \to I \to B \to \Gamma(\pi^{-1}\spec B,\O_X) \to 0 , $$
so that $\Gamma(\pi^{-1}\spec B,\O_X)=B/I$, and we have that $\pi^{-1}\spec B=\spec(B/I)\subseteq\spec B$ (viewing $X$ as a subspace of $Y$).

Conversely if $\I$ is a quasicoherent sheaf of ideals, so it fits into an exact sequence where $\F$ is quasicoherent on $Y$
$$ 0 \to \I \to \O_Y \to \F \to 0 . $$
For $\spec B\subseteq Y$ affine open let $\I(B)=\Gamma(\spec B,\I)$.
Now notice that for $\spec A,\spec B\subseteq Y$ affine open, $\spec A/\I(A)\subseteq\spec A$ and $\spec B/\I(B)\subseteq\spec B$
restrict to the same closed subscheme of the intersection $\spec A\cap\spec B$.
We may cover the intersection $\spec A\cap\spec B$ with open affine subschemes which are simultaneously distinguished in both $\spec A$ and
$\spec B$.

Now because $\I$ is quasicoherent then $\I|_{\spec A}=\tilde{\I(A)}$ and so $\I(\spec A_f)=\I(A)_f$.
Thus if $\spec A_f=\spec B_g\subseteq\spec A\cap\spec B$ we have $\I(A)_f=\I(B)_g$ and $A_f=B_g$.
Thus $\spec A_f/\I(A)_f=\spec B_g/\I(B_g)$, and because $\spec A_f/\I(A)_f=\spec A/\I(A)\cap\spec A_f\subseteq\spec A$, we see that we can glue
$\spec A/\I(A)$ together to get a closed subscheme $X$ of $Y$.

So we have proven the following.

\bprop

    $\pi\colon X\inj Y$ is a closed embedding iff $\pi$ is a closed topological embedding, $\pi^\sharp$ is surjective, and $\ker(\pi^\sharp)$
    is a quasicoherent sheaf of ideals.

    Conversely if $\I$ is a quasicoherent sheaf of ideals, then it defines a closed subscheme of $Y$.

\eprop

\bremark

    In fact we need not require that $\ker(\pi^\sharp)$ be quasicoherent, this can be proven from the first
    two assumptions (see homework 2).

\eremark

\bdefn

    Let $Z_i\subseteq X$ be (arbitrarily many) closed subschemes of $X$.
    These correspond to quasicoherent sheaves of ideals $\I_i\leq\O_X$.
    Their sum $\sum_i\I_i$ is a quasicoherent sheaf of ideals, and thus corresponds to a closed
    subscheme of $X$ which we call the {\emphcolor scheme-theoretic intersection} of $Z_i$.

    If there are only finitely many $Z_i$, then $\bigcap_i\I_i$ is a quasicoherent sheaf of ideals
    and thus corresponds to a closed subscheme of $X$ called the {\emphcolor scheme-theoretic union}
    of $Z_i$.

\edefn

\bexam

    Note that for any affine scheme $X=\spec A$, its closed subschemes are of the form $\spec(A/I)=\vof_A(I)$.
    Then the scheme theoretic intersection and union of closed subschemes are the usual affine
    intersection and union:
    $$ \bigcap_i\vof_A(I_i) = \vof\parens{\sum_iI_i},\qquad \vof_A(I)\cup\vof_A(J) = \vof_A(IJ) =
    \vof_A(I\cap J) . $$

\eexam

In order for these definitions to be valid we must show that the sum of quasicoherent sheaves of ideals
is quasicoherent, and similarly for the intersection.
This is clear as if $\I_i|_{\spec A}=\tilde{I_i}$ then
$\eval{\parens{\sum_i\I_i}}_{\spec A}=\tilde{\sum_iI_i}$.
And similarly the intersection of finitely many quasicoherent sheaves of ideals is quasicoherent
(because localization by $f\in A$ commutes with finite intersections of ideals).

Furthermore note that the underlying set of $\bigcap_iZ_i$ is the intersection of $Z_i$'s underlying sets,
and similar for union.
This is true affine-locally and thus everywhere.

We see that closed embeddings are monic.
Indeed suppose $\pi\colon Z\inj X$ is a closed embedding and $\phi_i\colon Y\to Z$ are morphisms with
$\pi\circ\phi_1=\pi\circ\phi_2$.
Then as topological maps $\phi=\phi_1=\phi_2$ (because $\pi$ is an injective topological map), and
$$ (\pi\circ\phi_i)^\sharp = \O_X \xtos{\pi^\sharp} \pi_*\O_Z \xtos{\pi_*\phi_i^\sharp} \pi_*\phi_*\O_Y . $$
And because $\pi^\sharp$ is surjective, $\pi_*\phi_1^\sharp=\pi_*\phi_2^\sharp$.
Because $\pi$ is a closed topological embedding we have that for $p\in Z$, the stalk of $\pi_*\phi_i^\sharp$
is just the stalk of $\phi_i^\sharp$ at $p$.
Thus $\phi_1^\sharp=\phi_2^\sharp$ as they agree at all stalks of $Z$.

So we have proven (the case for open embeddings is easy).

\bprop

    A closed or open embedding of schemes is monic.

\eprop

Let $Y$ be a scheme and $s\in\Gamma(Y,\O_Y)$ a global section.
This defines a presheaf of ideals, $(s)\leq\O_Y$ defined on $U\subseteq Y$ by
$\Gamma(U,(s))=(s|_U)\leq\Gamma(U,\O_Y)$ (it is the principal ideal generated by $s$'s restriction).
This is quasicoherent because for $a\in A$ and $f\in A$ we have $(a_f)=(a)_f\leq A_f$ and so
$(s)|_{\spec A}=\tilde{(s|_{\spec A})}$, and thus it is a quasicoherent sheaf.

\bdefn

    For a global section $s\in\Gamma(Y,\O_Y)$, the closed subscheme of $Y$ whose associated quasicoherent
    sheaf of ideals is $(s)$ is called the {\emphcolor vanishing scheme of $s$}, denoted $\vof(s)$.
    For a subset $S\subseteq\Gamma(Y,\O_Y)$ we define its vanishing scheme by
    $$ \vof(S) = \bigcap_{s\in S}\vof(s) . $$

\edefn

Affine-locally $\vof(s)$ has the form $\vof(s)|_{\spec B}=\vof_B(s|_{\spec B})$.
We may extend the above discussion to any ideal $I\leq\Gamma(Y,\O_Y)$ and define $\I$ similarly.
Then it is easy to see that $\vof(S)=\vof((S))$ where $(S)\leq\Gamma(Y,\O_Y)$ is the ideal generated by $S$.

\subsubsection{The scheme-theoretic image}

\bdefn

    Let $\iota\colon Z\inj Y$ be a closed subscheme, giving an exact sequence
    $0\to\I_{Z/Y}\to\O_Y\to\iota_*\O_Z\to0$.
    For a morphism $\pi\colon X\to Y$, we say that its {\emphcolor image lies in $Z$} if
    $\I_{Z/Y}\to\O_Y\to\pi_*\O_X$ is zero.
    That is (locally) if a function of $Y$ vanishes on $Z$, then it pulls back to the zero function on $X$.

    If $\pi$'s image lies in $Z_i$, then it also lies in their intersection (because
    $\I_i\to\O_Y\to\iota_*\O_Z$ are zero, so too must $\sum_i\I_i\to\O_Y\to\iota_*\O_Z$).
    So we define $\pi$'s {\emphcolor scheme-theoretic image} to be the smallest closed subscheme $Z$ in which
    $\pi$'s image lies (take the intersection of all such $Z$).

\edefn

When $Y=\spec B$ is affine then for $\pi\colon X\to Y$, its image is just $Z=\spec B/I$ where $I\leq B$ is
the kernel of $\pi^\sharp\colon B\to\Gamma(X,\O_X)$.
This is because $\I\to\O_Y\to\pi_*\O_X$ is zero iff it is zero on the global sections, iff $\I(B)$ is contained
in the kernel.

It would be nice to extend this affine-locally.
So for $\spec B\subseteq Y$ open affine let $\I(B)\leq B$ be the kernel of $B\to\Gamma(\spec B,\pi_*\O_X)$.
There is a natural map $\phi\colon\I(B)_g\to\I(B_g)$ for $g\in B$, induced by the map $\I(B)\to\I(B_g)$.
If this map is an isomorphism, then $\I$ forms a quasicoherent sheaf of ideals on $Y$.
So we want to determine when this map is an isomorphism.

$\phi$ is clearly injective: if $r\in B$ and $r/g^n$ is zero on $\pi^{-1}\spec B_g$ then because $g$
is invertible in $\spec B_g$, we have that $r$ must pull back to zero.
For surjectivity we need to determine when the following holds: given $r/g^n$ on $\spec B_g$ such that
it pulls back to zero on $\pi^{-1}\spec B_g$, is there some $m$ such that $rg^m$ pulls back to zero
on $\pi^{-1}\spec B$.

If $\pi^{-1}\spec B$ is reduced, then $\pi^\sharp(r)$ vanishes on $\pi^{-1}\dof g$, and $g$ vanishes on
$\vof(g)$, so that $\pi^\sharp(rg)$ vanishes on $\pi^{-1}\spec B=\pi^{-1}\dof(g)\cup\pi^{-1}\vof(g)$
(by considering that its value at each point is zero).

And if $\pi^{-1}\spec B=\spec A$ is affine, then let $r'=\pi^\sharp(r)$ and $g'=\pi^\sharp(g)$ in $A$.
Then if $r'=0$ on $\dof(g')=\pi^{-1}\dof g$, then there is an $m$ such that $r'(g')^m=0$ in $A$, so $rg^m$
pulls back to zero on $\spec B$.
And in general if $\pi^{-1}\spec B$ is quasicompact, then we can cover it with finitely many affine open
subsets such that $rg^{m_i}=0$ when pulling back, and take $m$ to be the maximum of $m_i$.

\bthrm

    Let $\pi\colon X\to Y$ be a morphism of schemes such that either $X$ is reduced or $\pi$ is quasicompact.
    Then the scheme-theoretic image of $X$ may be computed affine-locally: on $\spec A\subseteq Y$ it
    is $\spec A/I$ where $I$ is the kernel of $\pi^\sharp\colon A\to\Gamma(\spec A,\pi_*\O_X)$.

\ethrm

\bcoro

    The closure set-theoretic image of a morphism is contained in the underlying set of the scheme-theoretic
    image.
    They are equal under the conditions of the above theorem.

\ecoro

\bproof

    Let $\iota\colon Z\inj Y$ be a closed then $0\to\I_{Z/Y}\to\O_Y\to\iota_*\O_Z\to0$ is exact and thus
    for $p\notin Z$ we have that $(\iota_*\O_Z)_p=0$ so that $\I_{Z/Y,p}\cong\O_{Y,p}$.
    If the image of $\pi\colon X\to Y$ is contained in $Z$ then $\I_{Z/Y}\to\O_Y\to\pi_*\O_X$ is zero.
    Taking stalks this means that $\O_{Y,p}\to(\pi_*\O_X)_p$ is zero.
    Because these are rings this means $(\pi_*\O_X)_p$ is zero.
    If $p$ is in the image of $\pi$ then there would have to be an open $p\in U\subseteq Y$ such that
    $\pi^{-1}U\neq\emptyset$ and thus the colimit is not zero, a contradiction.
    So the set-theoretic image of $\pi$ is contained in the underlying set of the scheme theoretic image,
    and the scheme theoretic image is closed.

    Now let $Z^\isset$ be the set-theoretic image of $\pi$, and let $U=Y-Z^\isset$.
    Then $\pi^{-1}U=\emptyset$, and by the above theorem the scheme-theoretic image may be computed locally.
    So on $U$ we have that the scheme-theoretic image is empty.
    \qed

\eproof

Let $X^\isset\subseteq Y$ be a closed subset.
Consider the following three closed subschemes of $Y$.
\benum
    \item The smallest closed subscheme of $Y$ whose underlying set contains $X^\isset$.
    \item For an affine open $\spec B\subseteq Y$ consider the radical ideal $\I(B)=\I(X^\isset\cap\spec B)$,
    and glue together the $\spec(B/\I(B))$s to form a closed subscheme.
    It is easy to see that this is well-defined, as the ideals $I$ form a quasicoherent sheaf.
    \item Let $W$ be the scheme
    $$ W = \coprod_{p\in X^\isset}\spec\kappa(p) $$
    where $\kappa(p)=\O_{Y,p}/\m_{Y,p}$ is the residue field at $p$.
    There is a canonical map $\rho\colon W\to Y$ which sends the unique point of $\kappa(p)\in W$ to $p\in Y$.
    The coordinates of this are the morphisms $\spec\kappa(p)\to Y$ whose sheaf morphism corresponds to mapping
    a section to its value at $p$ (or zero if it lies in an open set disjoint from $p$).
    Then consider the closed subscheme which is the scheme-theoretic image of $\rho$.
\eenum
We claim all of these closed subschemes are equal.

(1) and (2) form the same closed subscheme, because if $X^\isset\subseteq Z$ then in $\spec B$ we must have
that $X^\isset\cap\spec B\subseteq Z\cap\spec B$, and $X^\isset\cap\spec B=\vof_B(\I(B))$ so that
$\I_{Z/Y}(B)\subseteq\I(B)$.
Thus $\I$ is the largest such quasicoherent sheaf and thus defines the smallest such closed subscheme.
Notice that this shows that the underlying set of this scheme is $X^\isset$.

For (3) notice that the stalk of $\rho_*\O_W$ at $p\in Y$ is $\kappa(p)$ when $p\in X^\isset$ and zero
otherwise.
Thus $\I_{Z/Y}\to\O_Y\to\rho_*\O_W$ is zero iff for every $p\in X^\isset$ the map
$\I_{Z/Y,p}\to\O_{Y,p}\to\kappa(p)$ is zero, iff $\I_{Z/Y,p}\subseteq\m_{Y,p}$ in $\O_{Y,p}$.
Because $\I_{Z/Y,p}$ is an ideal of $\O_{Y,p}$ this is iff it is proper.
Now $\I_{Z/Y,p}$ is only not proper when $p\in Z$, because it is all of $\O_{Y,p}$ only when
$(\iota_*\O_Z)_p=0$ which is iff $p\notin Z$.
And so the image of $\rho$ is contained in $Z$ iff $X^\isset\subseteq Z$.
That is, its image is the smallest closed subscheme containing $X^\isset$.

By (2) it is clear that this closed subscheme is reduced, as $B/\I(B)$ is.
This motivates the following definition.

\bdefn

    Let $X\subseteq Y$ be a closed sub{\it set\/} of the scheme $Y$.
    Then the closed subscheme constructed above (characterized by any of (1), (2), or (3)) is called the
    {\emphcolor reduced subscheme structure} on $X$ and is denoted $X^\red$.

    When $X=Y$ then $Y^\red$ is called the {\emphcolor reduction} of $Y$.

\edefn

Note that on the affine open $\spec B\subseteq Y$, $Y^\red$ corresponds to $\spec B/\nil(B)$.

